% ==============================================================================
% CH02_PYTHON_MASTERY.TEX — CHƯƠNG 2: KỸ NGHỆ PYTHON CHUYÊN SÂU CHO KỸ SƯ AI
% ==============================================================================
% Định dạng: Sách chuyên khảo kinh điển (Classic Academic & Engineering Monograph)
% Chuẩn mực: Không dùng hộp đóng khung 4 phía, toán học 5 bước, bài tập đa tầng
% ==============================================================================

\chapter{Kỹ Nghệ Python Chuyên Sâu Dành Cho Kỹ Sư AI}
\label{ch:python-mastery}

Python từ lâu đã trở thành ngôn ngữ thống trị trong toàn bộ hệ sinh thái Trí tuệ Nhân tạo và Dữ liệu lớn. Tuy nhiên, khoảng cách giữa một nhà nghiên cứu viết mã thử nghiệm trong Jupyter Notebook và một Kỹ sư AI xây dựng các hệ thống phục vụ suy luận (\textit{Inference Serving}) chịu tải hàng chục nghìn yêu cầu mỗi giây là vô cùng lớn. 

Một kỹ sư AI chuyên nghiệp không thể chỉ dừng lại ở cú pháp cơ bản; họ bắt buộc phải thấu hiểu tường minh cấu trúc nội vi của CPython, cơ chế cấp phát bộ nhớ tầng thấp, nguyên lý hoạt động của bộ thu gom rác thế hệ, các hạn chế và giải pháp vượt qua Khóa Thông dịch Toàn cục (GIL), mô hình lập trình bất đồng bộ với TaskGroup, kỹ thuật zero-copy memoryview, và làm chủ các kỹ thuật profiling tối ưu hóa hiệu năng vi mô.

---

\section{Cơ Chế Nội Vi Của CPython \& Cấu Trúc Bộ Nhớ Tầng Thấp}
\label{sec:cpython-internals}

\subsection{Bản Chất Thực Thi Của CPython: Từ Mã Nguồn Đến Máy Ảo}

Để hiểu tại sao mã Python lại có những đặc tính về hiệu năng và bộ nhớ như hiện nay, chúng ta phải bắt đầu từ vòng đời thực thi của một tập lệnh. CPython (bản cài đặt tham chiếu chuẩn bằng ngôn ngữ C) không trực tiếp biên dịch mã nguồn thành mã máy nhị phân của CPU như C hay Rust, mà trải qua một quy trình nhiều giai đoạn:

\begin{enumerate}
  \item \textbf{Phân tích từ vựng và ngữ pháp (Lexing \& Parsing)}: Trình thông dịch đọc tệp nguồn \code{.py}, phân rã thành luồng token, và xây dựng Cây Cú pháp Trừu tượng (\textit{Abstract Syntax Tree -- AST}).
  \item \textbf{Biên dịch thành Bytecode (Bytecode Compilation)}: AST được bộ biên dịch Python chuyển hóa thành một chuỗi các chỉ lệnh bytecode trung gian độc lập với phần cứng. Kết quả này được đóng gói trong đối tượng mã (\code{PyCodeObject}) và lưu tạm vào đĩa dưới dạng tệp nhị phân \code{.pyc} trong thư mục \code{__pycache__}.
  \item \textbf{Vòng lặp thông dịch của máy ảo (Evaluation Loop)}: Trình thông dịch CPython (chủ yếu là hàm \code{_PyEval_EvalFrameDefault} trong tệp \code{Python/ceval.c}) thực thi một vòng lặp vô tận, đọc từng lệnh bytecode, sử dụng cấu trúc máy ảo dựa trên ngăn xếp (\textit{Stack-based Virtual Machine}) để đẩy (\textit{push}) và lấy (\textit{pop}) dữ liệu từ ngăn xếp định giá.
\end{enumerate}

Vì mọi chỉ lệnh bytecode đều phải đi qua vòng lặp định giá trừu tượng của C và thực hiện kiểm tra kiểu dữ liệu động tại thời điểm chạy (\textit{Runtime Dynamic Dispatch}), tốc độ thực thi thuần túy của Python chậm hơn ngôn ngữ C từ $10$ đến $50$ lần cho các vòng lặp CPU. Do đó, kỹ sư AI cần nắm vững cơ chế đẩy các tác vụ tính toán nặng xuống các phần mở rộng C/C++ (như NumPy, PyTorch, vLLM Kernels) thông qua Python C-API.

\subsection{Cấu Trúc Cốt Lõi \code{PyObject} \& \code{PyVarObject}}

Trong CPython, tất cả các thực thể (số nguyên, số thực, chuỗi, hàm, lớp, và các mô hình học máy) đều là một con trỏ trỏ tới cấu trúc \code{PyObject} trên vùng nhớ Heap.

Về mặt mã nguồn C (trong \code{Include/object.h}), cấu trúc cơ bản của \code{PyObject} được định nghĩa như sau:

\begin{lstlisting}[language=C]
// Dinh nghia cot loi trong CPython Include/object.h
typedef struct _object {
    _PyObject_HEAD_EXTRA         // Con tro kep danh cho danh sach lien ket cua Cyclic GC
    Py_ssize_t ob_refcnt;        // Bo dem tham chieu (Reference Count: 8 bytes tren 64-bit)
    struct _typeobject *ob_type; // Con tro tro toi kieu du lieu cua doi tuong (8 bytes)
} PyObject;
\end{lstlisting}

Đối với các đối tượng có kích thước thay đổi được theo thời gian hoặc độ dài (như chuỗi ký tự, danh sách, bộ dữ liệu số nguyên độ dài tùy ý), CPython mở rộng cấu trúc trên thành \code{PyVarObject}:

\begin{lstlisting}[language=C]
typedef struct {
    PyObject ob_base;            // Phan dau PyObject (16 bytes)
    Py_ssize_t ob_size;          // So luong phan tu hien tai (8 bytes tren 64-bit)
} PyVarObject;
\end{lstlisting}

\begin{table}[htbp]
\centering
\small
\caption{Bảng phân tích bộ nhớ từng byte của các đối tượng nguyên thủy trong CPython 64-bit}
\label{tab:cpython-memory-layout}
\begin{tabular}{p{2.5cm}p{2.6cm}cp{6.5cm}}
\toprule
\textbf{Kiểu dữ liệu} & \textbf{Cấu trúc C tương ứng} & \textbf{Bộ nhớ tối thiểu} & \textbf{Chi tiết phân bổ từng byte} \\
\midrule
Số nguyên C (\code{int64\_t}) & Giá trị thuần & 8 bytes & 8 bytes giá trị nhị phân thuần túy trên thanh ghi hoặc stack. \\
Python \code{int} & \code{PyLongObject} & 28 bytes & 8B \code{ob_refcnt} + 8B \code{ob_type} + 8B \code{ob_size} + 4B \code{ob_digit[0]}. \\
Python \code{float} & \code{PyFloatObject} & 24 bytes & 8B \code{ob_refcnt} + 8B \code{ob_type} + 8B giá trị \code{double} (IEEE 754). \\
Python \code{str} (1 ký tự ASCII) & \code{PyASCIIObject} & 50 bytes & 16B \code{PyObject} + 8B \code{length} + 8B \code{hash} + 8B \code{state} + 8B \code{wchar} + 2B (char + null). \\
Python \code{list} (rỗng) & \code{PyListObject} & 56 bytes & 24B \code{PyVarObject} + 8B con trỏ mảng \code{**ob_item} + 8B \code{allocated} + 16B GC header. \\
Python \code{dict} (rỗng) & \code{PyDictObject} & 64 bytes & Cấu trúc Compact Dict với mảng băm và mảng chỉ mục rỗng. \\
\bottomrule
\end{tabular}
\end{table}

\begin{definition}[Chi Phí Bộ Nhớ Phụ Trội -- PyObject Memory Overhead]
Một số thực IEEE 754 64-bit trong C chỉ tiêu tốn đúng $8\text{ bytes}$. Tuy nhiên, trong Python, một đối tượng \code{float} tiêu tốn tối thiểu $24\text{ bytes}$ (gấp 3 lần). Nghiêm trọng hơn, một danh sách Python (\code{list[float]}) chứa $N$ số thực không chỉ tốn $N \times 24\text{ bytes}$ cho các đối tượng số thực mà còn tốn thêm $N \times 8\text{ bytes}$ cho mảng con trỏ nội vi trỏ tới các đối tượng đó, cộng thêm bộ nhớ đệm dự phòng (\textit{Over-allocation}). Tổng chi phí lên tới $32\text{ bytes}$ cho mỗi phần tử -- gấp $4$ lần mảng C thuần túy!
\end{definition}

\subsection{Cơ Chế Cấp Phát Danh Sách \& Bộ Từ Điển Compact Dict}

\paragraph{Chiến lược cấp phát dư thừa của List (\code{PyListObject})}:
Khi bạn liên tục thêm phần tử vào danh sách bằng phương thức \code{list.append()}, CPython không thực hiện gọi hàm \code{realloc} của hệ điều hành cho từng phần tử một. Thay vào đó, thuật toán trong hàm \code{list_resize()} áp dụng công thức tăng trưởng kích thước tỷ lệ thuận:

\begin{equation}
\text{allocated} = \text{newsize} + (\text{newsize} \gg 3) + (\text{newsize} < 9 \; ? \; 3 : 6)
\label{eq:list-overallocation}
\end{equation}

Công thức này đảm bảo chi phí thời gian cho phép thêm phần tử khấu hao (\textit{Amortized Time Complexity}) đạt $\mathcal{O}(1)$. Tuy nhiên, điều này cũng đồng nghĩa với việc danh sách có thể đang giữ một vùng nhớ đệm dư thừa từ $12.5\%$ đến $25\%$ không sử dụng. Đối với các pipeline xử lý dữ liệu huấn luyện hàng trăm triệu mẫu, việc chuyển đổi danh sách thành các mảng NumPy dạng liên tục (\code{C-contiguous}) là điều kiện sống còn để bảo vệ bộ nhớ RAM không bị cạn kiệt.

\paragraph{Kiến trúc bộ từ điển Compact Dict (PEP 468/412)}:
Trong các phiên bản Python cũ (trước 3.6), từ điển được cài đặt dưới dạng một bảng băm thưa thớt (\textit{Sparse Hash Table}), trong đó mỗi ô chứa đồng thời mã băm (\code{hash}), con trỏ khóa (\code{key}), và con trỏ giá trị (\code{value}), tiêu tốn $24\text{ bytes}$ cho mỗi ô bất kể ô đó có dữ liệu hay trống rỗng.

Từ Python 3.6+, CPython chuyển sang kiến trúc \textbf{Compact Dict}:
\begin{itemize}
  \item Mảng chỉ mục (\code{indices}): Một mảng nhỏ các số nguyên (1 byte kiểu \code{int8_t} nếu kích thước nhỏ) đại diện cho bảng băm thưa.
  \item Mảng thực thể (\code{entries}): Một mảng dày đặc (\textit{dense array}) lưu trữ liên tục các phần tử \code{[hash, key, value]} theo đúng thứ tự chèn.
\end{itemize}
Kiến trúc này giúp giảm từ $30\%$ đến $40\%$ dung lượng RAM của các đối tượng từ điển, đồng thời giữ nguyên thứ tự chèn phần tử (\textit{Insertion-ordered dictionary}).

\subsection{Thuật Toán Đếm Tham Chiếu \& Bộ Thu Gom Rác Thế Hệ (Cyclic GC)}

CPython quản trị vòng đời bộ nhớ thông qua cơ chế lai hai tầng:

\paragraph{Tầng 1: Đếm tham chiếu (Reference Counting)}:
Trường \code{ob_refcnt} đếm số lượng tham chiếu đang trỏ tới đối tượng. Khi một tham chiếu được gán hoặc truyền vào hàm, \code{ob_refcnt} tăng $1$. Khi biến ra khỏi phạm vi (\textit{scope}) hoặc bị gán lại, \code{ob_refcnt} giảm $1$. Ngay khoảnh khắc:
\begin{equation}
\text{ob\_refcnt} == 0
\end{equation}
CPython lập tức thu hồi vùng nhớ của đối tượng về bộ phân bổ bộ nhớ nội vi (\code{PyMalloc}) mà không cần tạm dừng luồng thực thi của ứng dụng.

\paragraph{Tầng 2: Thu gom rác thế hệ (Generational Cyclic Garbage Collector)}:
Điểm yếu của cơ chế đếm tham chiếu là hiện tượng \textbf{tham chiếu vòng} (\textit{Reference Cycles}): Đối tượng A tham chiếu đến B, và B tham chiếu ngược lại đến A. Khi cả hai biến đại diện bị xóa, bộ đếm tham chiếu của chúng vẫn duy trì ở mức $1$, khiến chúng không bao giờ được giải phóng.

Bộ thu gom rác của Python giải quyết vấn đề này bằng cách quét các đối tượng chứa (\textit{Container Objects}) qua 3 thế hệ:
\begin{itemize}
  \item \textbf{Thế hệ 0 (Generation 0)}: Chứa tất cả các đối tượng vừa mới sinh ra. Tần suất quét rất cao khi số lượng cấp phát mới vượt quá ngưỡng (\textit{Threshold}).
  \item \textbf{Thế hệ 1 (Generation 1)}: Chứa các đối tượng sống sót qua lần quét của Thế hệ 0.
  \item \textbf{Thế hệ 2 (Generation 2)}: Chứa các đối tượng sống sót qua nhiều chu kỳ quét (thường là các đối tượng tĩnh tồn tại suốt vòng đời ứng dụng). Tần suất quét là thấp nhất.
\end{itemize}

Trong các đồ thị tính toán của PyTorch (\textit{Autograd Graph}) hoặc đồ thị thực thi trạng thái Agent (như LangGraph), các nút con thường xuyên giữ tham chiếu trỏ ngược về nút cha hoặc context của đồ thị. Nếu kỹ sư cấu hình không cẩn thận hoặc lưu trữ các closure bắt giữ biến cục bộ, Cyclic GC có thể gây ra hiện tượng giật cục (\textit{GC Pause}) đột ngột kéo dài hàng trăm mili-giây trên môi trường sản xuất.

\subsection{Khóa Thông Dịch Toàn Cục (GIL) \& Free-Threaded Python}

Khóa Thông dịch Toàn cục (\textit{Global Interpreter Lock -- GIL}) là một khóa Mutex bảo vệ toàn bộ cấu trúc nội vi của CPython, đảm bảo rằng tại bất kỳ thời điểm nào chỉ có duy nhất \textbf{một luồng thực thi mã bytecode Python} trên CPU, ngay cả khi máy chủ sở hữu $64$ hay $128$ lõi vật lý.

\paragraph{Tại sao GIL lại tồn tại?}
Mục đích cốt lõi của GIL là ngăn chặn xung đột dữ liệu (\textit{Race Condition}) trên trường \code{ob_refcnt}. Nếu hai luồng cùng lúc tăng hoặc giảm \code{ob_refcnt} mà không có cơ chế đồng bộ, lỗi phân bổ bộ nhớ kép (\textit{Double Free}) hoặc hỏng dữ liệu sẽ xảy ra ngay lập tức. Việc áp dụng khóa riêng lẻ cho từng đối tượng sẽ khiến hiệu năng của chương trình đơn luồng suy giảm nghiêm trọng do chi phí khóa đắt đỏ.

\paragraph{Cơ chế giải phóng GIL trong các thư viện AI}:
Các thư viện hiệu năng cao như NumPy, SciPy, PyTorch, TensorRT và vLLM giải quyết triệt để vấn đề này bằng cách giải phóng GIL khi thực hiện các phép toán nặng bằng C/C++:

\begin{lstlisting}[language=C]
// Minh hoa cach C-Extension giai phong GIL trong tinh toan song song
Py_BEGIN_ALLOW_THREADS
// Doan ma C thuc thi song song tren toan bo loi CPU bang OpenMP hoac CUDA
c_matrix_multiply_parallel(matrix_a, matrix_b, result, num_threads);
Py_END_ALLOW_THREADS
\end{lstlisting}

Nhờ macro \code{Py_BEGIN_ALLOW_THREADS}, các luồng Python khác (chẳng hạn như luồng nhận request HTTP hoặc luồng logging) hoàn toàn có thể chạy song song trong khi GPU hoặc CPU C-extension đang thực hiện nhân ma trận.

\paragraph{Xu thế Free-Threaded Python (PEP 703)}:
Bắt đầu từ Python 3.13, một chế độ thử nghiệm không có GIL (\textit{Free-threaded build}) đã được tích hợp. Để thay thế GIL, CPython sử dụng kỹ thuật đếm tham chiếu có bảo vệ (\textit{Biased Reference Counting}) và bộ nhớ phân trang không khóa (\textit{Mimalloc}), mở ra kỷ nguyên thực thi đa luồng song song thực sự trên CPU cho các ứng dụng Python thuần túy trong tương lai.

---

\section{Hệ Thống Kiểu Dữ Liệu Bất Biến, Metaprogramming \& Pydantic V2}
\label{sec:type-system-pydantic}

\subsection{Dynamic Typing vs Strict Typing Trong Production AI}

Trong phát triển phần mềm truyền thống, Python được yêu thích vì cú pháp linh hoạt và kiểu dữ liệu động (\textit{Duck Typing}). Tuy nhiên, trong các hệ thống AI hiện đại, các mô hình ngôn ngữ lớn (LLM) là các thực thể mang tính ngẫu nhiên (\textit{Stochastic Engines}). Khi một Agent yêu cầu LLM trích xuất thông tin dưới dạng JSON, mô hình có thể trả về một trường số dưới dạng chuỗi (\code{"port": "8080"}), thiếu các trường bắt buộc, hoặc sinh ra các thuộc tính hoàn toàn bất ngờ.

Nếu tầng backend xử lý trực tiếp đối tượng từ điển (\code{dict}) mà không có lớp thẩm định kiểu dữ liệu nghiêm ngặt, lỗi \code{KeyError} hoặc \code{TypeError} sẽ bùng phát ở sâu trong luồng nghiệp vụ gây sụp đổ hệ thống. Do đó, việc xây dựng một hệ thống kiểu dữ liệu bất biến và có cấu trúc (\textit{Strict Typed Schema}) là yêu cầu bắt buộc đối với một Kỹ sư AI.

\subsection{Kỹ Thuật Structural Subtyping Với \code{typing.Protocol}}

Thay vì sử dụng thừa kế danh nghĩa (\textit{Nominal Subtyping}) truyền thống với các lớp cơ sở trừu tượng (\code{abc.ABC}), Python hiện đại hỗ trợ \textbf{Structural Subtyping} (Duck Typing tĩnh) thông qua \code{typing.Protocol} (PEP 544). Điều này cho phép chúng ta áp dụng triệt để nguyên lý Nghịch đảo Phụ thuộc (\textit{Dependency Inversion Principle}):

\begin{lstlisting}[language=Python]
from typing import Protocol, Sequence, runtime_checkable

@runtime_checkable
class VectorIndex(Protocol):
    """Giao dien truu tuong cho kho luu tru vector khong rang buoc nha cung cap."""
    
    async def add_vectors(self, ids: Sequence[str], vectors: Sequence[Sequence[float]]) -> None:
        ...
        
    async def search_nearest(
        self, query_vector: Sequence[float], top_k: int
    ) -> Sequence[tuple[str, float]]:
        ...
\end{lstlisting}

Bất kỳ lớp nào cài đặt đầy đủ hai phương thức \code{add_vectors} và \code{search_nearest} đều tự động tương thích với giao diện \code{VectorIndex} mà không cần phải tường minh kế thừa từ một lớp cha nào. Điều này giúp mã nguồn của hệ thống tách rời hoàn toàn khỏi các thư viện cụ thể như Qdrant, Milvus hay ChromaDB.

\subsection{Cơ Chế Descriptors \& Metaprogramming Trong Framework AI}

Cơ chế Descriptors là nền tảng cốt lõi giúp các framework như PyTorch (\code{torch.nn.Module}), Pydantic, và SQLAlchemy vận hành ma thuật quản lý thuộc tính của chúng.

Một đối tượng là một Descriptor nếu nó định nghĩa ít nhất một trong các phương thức: \code{__get__}, \code{__set__}, hoặc \code{__delete__}. Khi một thuộc tính của lớp được truy cập:
\begin{equation}
\text{instance.attribute} \implies \text{type(instance).\_\_dict\_\_['attribute'].\_\_get\_\_(instance, type(instance))}
\end{equation}

\begin{lstlisting}[language=Python]
class BoundedTensorDescriptor:
    """Descriptor quan ly kiem tra bien va gia tri hop le cua Tensor."""
    
    def __init__(self, name: str, min_val: float = 0.0, max_val: float = 1.0):
        self.name = name
        self.min_val = min_val
        self.max_val = max_val
        self._storage_name = f"_{name}"

    def __get__(self, instance, owner):
        if instance is None:
            return self
        return getattr(instance, self._storage_name, None)

    def __set__(self, instance, value: float):
        if not (self.min_val <= value <= self.max_val):
            raise ValueError(f"{self.name} phai nam trong khoang [{self.min_val}, {self.max_val}]. Nhan duoc: {value}")
        setattr(instance, self._storage_name, value)
\end{lstlisting}

\subsection{Tối Ưu Hóa Bộ Nhớ Cực Hạn Với \code{\_\_slots\_\_}}

Mỗi thể hiện (\textit{instance}) của một lớp Python thông thường lưu trữ các thuộc tính của nó trong một từ điển nội vi \code{__dict__}. Mỗi từ điển này tiêu tốn tối thiểu $152\text{ bytes}$ bộ nhớ phụ trội. Khi bạn cần lưu trữ hàng triệu thực thể đồ thị tri thức (\textit{Knowledge Graph Entities}) hoặc các đoạn văn bản trong bộ nhớ đệm RAM, chi phí này trở thành thảm họa.

Bằng cách khai báo \code{__slots__}, CPython sẽ loại bỏ hoàn toàn \code{__dict__} và thay thế bằng một mảng con trỏ mảng tĩnh tương tự như một \code{struct} trong C:

\begin{lstlisting}[language=Python]
class ChunkNodeWithSlots:
    """Lop bieu dien chunk van ban toi uu hoa bo nho voi __slots__."""
    __slots__ = ("chunk_id", "doc_id", "embedding", "token_count")
    
    def __init__(self, chunk_id: str, doc_id: str, embedding: list[float], token_count: int):
        self.chunk_id = chunk_id
        self.doc_id = doc_id
        self.embedding = embedding
        self.token_count = token_count
\end{lstlisting}

Việc sử dụng \code{__slots__} giảm kích thước bộ nhớ của mỗi đối tượng từ khoảng $200\text{ bytes}$ xuống chỉ còn $56\text{ bytes}$ (tiết kiệm gần $75\%$ dung lượng RAM) và tăng tốc độ đọc ghi thuộc tính lên $20\%$.

\subsection{Kỹ Nghệ Schema \& Structured Output Với Pydantic V2}

Pydantic V2 được viết lại hoàn toàn với nhân thực thi bằng ngôn ngữ \textbf{Rust} (\code{pydantic-core}). Quá trình phân tích cú pháp JSON và thẩm định dữ liệu diễn ra hoàn toàn ở tầng C/Rust tốc độ cao mà không phải tạo các đối tượng Python trung gian.

\begin{lstlisting}[language=Python]
from pydantic import BaseModel, Field, field_validator, model_validator
from typing import Literal, Union
from datetime import datetime

class ToolCallTelemetry(BaseModel):
    """Schema bat bien xac thuc ket qua thuc thi cong cu cua LLM."""
    tool_name: str = Field(..., description="Ten cong cu duoc goi")
    execution_latency_ms: float = Field(..., ge=0.0)
    status: Literal["SUCCESS", "FAILED", "TIMEOUT"]
    output_tokens: int = Field(default=0, ge=0)
    created_at: datetime = Field(default_factory=datetime.utcnow)

    @field_validator("tool_name")
    @classmethod
    def validate_tool_name_safe(cls, v: str) -> str:
        if not v.replace("_", "").isalnum():
            raise ValueError(f"Ten cong cu chi duoc phep chua ky tu chu, so va gach duoi: {v}")
        return v
\end{lstlisting}

Pydantic V2 cho phép chúng ta định nghĩa các cấu trúc phức tạp như phân biệt kiểu động qua Discriminator Unions (\code{Field(discriminator='action_type')}), giúp kiểm soát chặt chẽ hàng nghìn cấu trúc Function Calling trong các hệ thống Multi-Agent lớn.

---

\section{Lập Trình Bất Đồng Bộ Nâng Cao Với AsyncIO \& Structured Concurrency}
\label{sec:asyncio-mastery}

\subsection{Bản Chất Vòng Lặp Sự Kiện (Event Loop) \& Multiplexing}

Khác với mô hình Đa tiến trình (\textit{Multiprocessing}) hoặc Đa luồng của Hệ điều hành (\textit{OS Threads}), \code{asyncio} hoạt động dựa trên cơ chế \textbf{Đa nhiệm Hợp tác (Cooperative Multitasking)} trên một luồng CPU duy nhất.

Cốt lõi của \code{asyncio} là vòng lặp sự kiện (\textit{Event Loop}) kết hợp với cơ chế I/O Multiplexing tầng nhân Linux (\code{epoll}):
\begin{itemize}
  \item Khi một Coroutine gửi một truy vấn mạng (ví dụ: gọi API suy luận tới mô hình LLM), nó đăng ký Socket File Descriptor với bộ kiểm soát \code{epoll} và thực thi lệnh \code{await}.
  \item Thao tác này ngay lập tức nhường quyền điều khiển CPU lại cho Event Loop.
  \item Trong thời gian chờ đợi phản hồi từ mạng (thường mất từ $50\text{ms}$ đến hàng giây đối với LLM), Event Loop tiếp tục xử lý hàng nghìn Coroutine khác mà không hề bị phong tỏa.
  \item Khi gói tin phản hồi đến card mạng vật lý, nhân Linux gửi tín hiệu đánh thức thông qua ngắt \code{epoll}, Event Loop tiếp tục thực thi Coroutine từ điểm đã tạm dừng.
\end{itemize}

\subsection{Hạn Chế Của \code{asyncio.gather} \& Kỷ Nguyên \code{TaskGroup}}

Trong Python phiên bản cũ, lập trình viên thường sử dụng lệnh \code{asyncio.gather(*tasks)} để chạy đồng thời nhiều tác vụ. Tuy nhiên, cách tiếp cận này ẩn chứa các nguy cơ nghiêm trọng về độ tin cậy trong các hệ thống sản xuất:
\begin{itemize}
  \item \textbf{Nguy cơ Tác vụ Mồ côi (Orphaned Tasks)}: Nếu một tác vụ trong \code{gather} gặp ngoại lệ nghiêm trọng và dừng lại, các tác vụ còn lại vẫn âm thầm chạy ngầm trong bộ nhớ mà không có cơ chế hủy tự động, tiêu tốn tài nguyên và tạo ra trạng thái dữ liệu không nhất quán.
  \item \textbf{Nguy cơ Nuốt Ngoại Lệ (Exception Swallowing)}: Việc xử lý nhiều ngoại lệ xảy ra đồng thời rất phức tạp và dễ gây rò rỉ.
\end{itemize}

Từ Python 3.11+, triết lý \textbf{Structured Concurrency} được giới thiệu chính thức thông qua cấu trúc \code{TaskGroup} của module \code{asyncio}:

\begin{lstlisting}[language=Python]
import asyncio

async def fetch_embedding_shard(shard_id: int) -> list[float]:
    await asyncio.sleep(0.01 * shard_id)
    if shard_id == 3:
        raise ConnectionResetError("Shard 3 bi mat ket noi mang")
    return [0.1 * shard_id] * 4

async def query_all_shards():
    """Structured Concurrency: Neu 1 shard loi, toan bo cac shard con lai bi huy ngay lap tuc."""
    results = {}
    try:
        async with asyncio.TaskGroup() as tg:
            for shard_id in range(5):
                task = tg.create_task(fetch_embedding_shard(shard_id))
                results[shard_id] = task
    except* ConnectionResetError as eg:
        # Xu ly ngoai le theo nhom (Exception Groups - PEP 654)
        print(f"[ERROR CAUGHT] Phat hien loi ket noi: {eg.exceptions}")
\end{lstlisting}

Với \code{TaskGroup}, khi khối \code{async with} kết thúc, toàn bộ các tác vụ con được đảm bảo đã hoàn tất hoặc bị hủy bỏ triệt để. Nếu một tác vụ ném ngoại lệ, tất cả các tác vụ còn lại trong nhóm sẽ ngay lập tức nhận tín hiệu \code{cancel()}, triệt tiêu hoàn toàn nguy cơ tác vụ mồ côi.

\subsection{Kiểm Soát Tải Với Async Semaphore \& Token Bucket Rate Limiter}

Trong hạ tầng phục vụ AI, việc kiểm soát tốc độ gửi truy vấn (\textit{Rate Limiting}) là bắt buộc để không làm sụp đổ các cụm máy chủ suy luận vLLM hoặc bị khóa tài khoản do vượt hạn mức nhà cung cấp.

Thuật toán chuẩn mực nhất để quản lý lưu lượng là \textbf{Token Bucket (Thùng chứa thẻ)}:
\begin{itemize}
  \item Thùng có dung lượng tối đa $B$ thẻ (Tokens).
  \item Thẻ được bổ sung vào thùng với tốc độ không đổi $r$ thẻ mỗi giây.
  \item Mỗi truy vấn cần tiêu tốn $k$ thẻ (tương ứng với số lượng token đầu vào hoặc số yêu cầu HTTP). Nếu thùng còn đủ thẻ, truy vấn được phép đi qua; nếu không, truy vấn phải chờ cho đến khi thẻ được tích lũy đủ.
\end{itemize}

\begin{figure}[htbp]
\centering
\small
\begin{verbatim}
       Tốc độ nạp thẻ r (tokens/giây)
                 |
                 v
        +------------------+
        |   Token Bucket   |  Dung lượng tối đa: B
        |   o  o  o  o  o  |  Số thẻ hiện tại: tokens(t)
        +------------------+
                 |
                 v (Cần k thẻ cho mỗi Request)
        [Kiểm tra khả dụng]
           /            \
     (Đủ thẻ)        (Thiếu thẻ)
        /                \
       v                  v
[Cho phép gửi ngay]   [Chờ t_wait giây]
\end{verbatim}
\caption{Mô hình giải thuật Token Bucket kiểm soát thông lượng cho AI Gateway}
\label{fig:token-bucket-flow}
\end{figure}

Công thức tính toán thời gian chờ đợi chính xác khi lượng thẻ trong thùng không đủ $k$ thẻ:
\begin{equation}
T_{\text{wait}} = \frac{k - \text{tokens}_{\text{current}}}{r}
\label{eq:token-bucket-wait}
\end{equation}

---

\section{Kỹ Nghệ Zero-Copy, Buffer Protocol \& Tối Ưu Dòng Dữ Liệu}
\label{sec:zero-copy-memoryview}

\subsection{Chi Phí Ẩn Của Thao Tác Sao Chép Dữ Liệu (Memory Copy)}

Khi xây dựng các dịch vụ phục vụ AI xử lý âm thanh, hình ảnh hoặc các ma trận trọng số mô hình lớn, một trong những điểm nghẽn hiệu năng lớn nhất là \textbf{sao chép bộ nhớ ngầm định}.

Mỗi khi bạn thực hiện cắt một chuỗi bytes (\code{chunk = data[start:end]}), Python tạo ra một đối tượng \code{bytes} hoàn toàn mới và sao chép toàn bộ khối byte từ vùng nhớ nguồn sang vùng nhớ đích. Nếu bạn xử lý một tệp âm thanh WAV $100\text{ MB}$ và chia thành $1.000$ đoạn nhỏ, bạn đã buộc CPU phải sao chép $100\text{ MB}$ dữ liệu qua lại trong RAM, gây nghẽn băng thông bộ nhớ và kích hoạt bộ thu gom rác dọn dẹp hàng nghìn đối tượng thừa.

\subsection{Python Buffer Protocol \& Lớp \code{memoryview}}

Python cung cấp giao thức \textbf{Buffer Protocol} ở tầng C-API, cho phép các đối tượng chia sẻ chung một bộ nhớ đệm nhị phân liên tục mà không cần sao chép. Thông qua đối tượng \code{memoryview}, kỹ sư có thể cắt mảng, trỏ vào các phân vùng dữ liệu với chi phí thời gian $\mathcal{O}(1)$ và tiêu tốn đúng $0\text{ bytes}$ bộ nhớ cấp phát mới:

\begin{lstlisting}[language=Python]
import time

def demo_zero_copy_performance():
    # Tao khoi du lieu nhi phan 50 MB
    raw_data = bytearray(50 * 1024 * 1024)
    
    # Cach 1: Copy truyen thong (Tieu ton 50MB RAM moi lan cat)
    t0 = time.perf_counter()
    chunk1 = raw_data[1024:2048] # Tao ban sao moi tren RAM
    t_copy = (time.perf_counter() - t0) * 1000.0
    
    # Cach 2: Zero-copy Memoryview (Khong copy bat ky byte nao)
    t1 = time.perf_counter()
    view = memoryview(raw_data)
    chunk2 = view[1024:2048] # Chi tao con tro tham chieu truc tiep
    t_view = (time.perf_counter() - t1) * 1000.0
    
    print(f"Thoi gian Slice truyen thong: {t_copy:.4f}ms | Zero-copy View: {t_view:.4f}ms")
\end{lstlisting}

Kỹ thuật Zero-copy này là nguyên lý căn bản đằng sau cách mà PyTorch chia sẻ bộ nhớ giữa Tensor CPU và NumPy ndarray (\code{torch.from_numpy(arr)}), cũng như cách các máy chủ Triton truyền tải payload tensor qua IPC Shared Memory (\code{/dev/shm}).

---

\section{Kỹ Thuật Tracing, Profiling \& Tối Ưu Hóa Hiệu Năng Vi Mô}
\label{sec:profiling-tuning}

Để tối ưu hóa một hệ thống AI, phương châm sống còn của kỹ sư là: \textbf{Không bao giờ đoán mò điểm nghẽn; luôn luôn đo lường bằng Profiler chuyên dụng.}

\subsection{Profiling CPU Với \code{cProfile} \& Profiler Không Xâm Lấn \code{py-spy}}

\begin{itemize}
  \item \textbf{\code{cProfile}}: Profiler tích hợp sẵn của Python, ghi nhận số lần gọi hàm (\code{ncalls}) và tổng thời gian thực thi tích lũy (\code{cumtime}). Thích hợp cho việc phân tích trong môi trường kiểm thử cục bộ:
  \begin{lstlisting}[language=bash]
  python -m cProfile -s cumtime service_server.py
  \end{lstlisting}
  \item \textbf{\code{py-spy}}: Profiler lấy mẫu không xâm lấn (\textit{Sampling Profiler}) được viết bằng Rust. Điểm ưu việt tuyệt đối của \code{py-spy} là nó có thể đính kèm vào một tiến trình Python đang chạy trên máy chủ sản xuất (\textit{Production}) mà không cần khởi động lại dịch vụ và không làm suy giảm tốc độ xử lý:
  \begin{lstlisting}[language=bash]
  # Ghi lai bieu do ngon lua (Flamegraph) truc tiep tu tien trinh AI dang chay
  py-spy record -o flamegraph.svg --pid 12345 --duration 30
  \end{lstlisting}
\end{itemize}

\subsection{Memory Profiling \& Truy Vết Rò Rỉ RAM Với \code{tracemalloc}}

Trong các hệ thống Agent phục vụ liên tục trong nhiều tháng, hiện tượng rò rỉ bộ nhớ (\textit{Memory Leak}) tích tụ dần có thể dẫn đến việc hệ điều hành tiêu diệt tiến trình bằng Out-Of-Memory (OOM Killer).

Module \code{tracemalloc} cho phép ghi nhận vết cấp phát bộ nhớ của từng dòng mã nguồn:

\begin{lstlisting}[language=Python]
import tracemalloc

def trace_memory_leak():
    tracemalloc.start()
    snapshot1 = tracemalloc.take_snapshot()
    
    # Thuc thi tac vu can kiem tra ro ri
    leaky_cache = []
    for i in range(10000):
        leaky_cache.append(f"token_id_data_leak_{i}")
        
    snapshot2 = tracemalloc.take_snapshot()
    top_stats = snapshot2.compare_to(snapshot1, 'lineno')
    
    print("[TOP 3 DIEM CAP PHAT BO NHO NHIEU NHAT]")
    for stat in top_stats[:3]:
        print(stat)
\end{lstlisting}

---

\section{Ví Dụ Giải Mẫu Tính Toán Số Học Bằng Tay}
\label{sec:ch02-numerical-walkthrough}

\begin{example}[Chi Phí Lưu Trữ 1 Triệu Vector Embeddings: List Thuần vs Mảng NumPy]
\label{ex:memory-vectors}
Giả sử hệ thống cần lưu trữ trong bộ nhớ RAM $N = 1.000.000$ vector embeddings phục vụ tìm kiếm tương đồng tức thì. Mỗi vector có số chiều $d = 1.536$ (chuẩn của mô hình OpenAI \code{text-embedding-3-small}).

\paragraph{Trường hợp 1: Sử dụng cấu trúc danh sách lồng nhau chuẩn Python (\code{list[list[float]]})}:
\begin{enumerate}
  \item Mỗi số thực trong Python là một đối tượng \code{PyFloatObject} tốn $24\text{ bytes}$.
  \item Mảng con trỏ của danh sách trong cùng cần $8\text{ bytes}$ cho mỗi con trỏ trỏ tới đối tượng \code{float}.
  \item Bộ nhớ cho 1 vector $1.536$ chiều:
  \begin{align*}
  M_{\text{single-vec}} &= d \times (24\text{ B (PyFloatObject)} + 8\text{ B (Pointer)}) + 56\text{ B (List Header)} \\
  &= 1.536 \times 32 + 56 = 49.208\text{ bytes} \approx 48.05\text{ KiB}
  \end{align*}
  \item Danh sách ngoài cùng chứa $1.000.000$ con trỏ tới các danh sách con:
  \begin{equation*}
  M_{\text{outer-list}} \approx 10^6 \times 8\text{ bytes} + 56\text{ bytes} \approx 8.000.056\text{ bytes} \approx 7.63\text{ MiB}
  \end{equation*}
  \item Tổng dung lượng RAM cần thiết:
  \begin{align*}
  M_{\text{Total, Python List}} &= 10^6 \times 49.208\text{ bytes} + 8\times 10^6\text{ bytes} \\
  &= 49.216 \times 10^6\text{ bytes} \approx 45.84\text{ GiB}!
  \end{align*}
  Hệ quả: Bạn cần một máy chủ Cloud cỡ lớn với $64\text{ GB}$ RAM chỉ để chứa dữ liệu thô!
\end{enumerate}

\paragraph{Trường hợp 2: Sử dụng mảng nhị phân liên tục NumPy Float32 (\code{numpy.ndarray})}:
\begin{enumerate}
  \item Không có vỏ bọc \code{PyObject} cho từng số riêng lẻ. Toàn bộ $1.000.000 \times 1.536$ phần tử được đặt sát nhau trên một khối bộ nhớ phẳng liên tục.
  \item Mỗi số thực chuẩn Float32 tốn đúng $4\text{ bytes}$ ($32\text{ bits}$).
  \item Tổng dung lượng bộ nhớ cho ma trận:
  \begin{equation*}
  M_{\text{Total, NumPy FP32}} = 10^6 \times 1.536 \times 4\text{ bytes} + 128\text{ B (Header)} = 6.144.000.128\text{ bytes} \approx 5.72\text{ GiB}!
  \end{equation*}
\end{enumerate}

\paragraph{Trường hợp 3: Áp dụng Lượng tử hóa Int8 (\code{int8} Quantization)}:
Nếu chuyển đổi sang biểu diễn số nguyên 8-bit ($1\text{ byte}$ cho mỗi chiều):
\begin{equation*}
M_{\text{Total, Int8}} = 10^6 \times 1.536 \times 1\text{ byte} = 1.536.000.000\text{ bytes} \approx 1.43\text{ GiB}!
\end{equation*}

\paragraph{Kết luận}: Việc thấu hiểu cấu trúc bộ nhớ tầng thấp giúp tiết kiệm \textbf{từ 8 đến 32 lần dung lượng RAM} ($1.43\text{ GiB}$ so với $45.84\text{ GiB}$), biến một hệ thống cồng kềnh thành một kiến trúc siêu nhẹ có thể vận hành trên một máy chủ giá rẻ.
\end{example}

\begin{example}[Tính Toán Thời Gian Chờ Của Thuật Toán Token Bucket Khi Bị Đột Biến Tải]
\label{ex:token-bucket-math}
Một API Gateway giới hạn lưu lượng gọi LLM với cấu hình Token Bucket:
\begin{itemize}
  \item Dung lượng tối đa của thùng: $B = 100$ thẻ.
  \item Tốc độ bổ sung thẻ: $r = 10$ thẻ/giây.
  \item Tại thời điểm $t = 0$, thùng chứa đầy $100$ thẻ.
  \item Đột ngột xuất hiện một loạt 3 yêu cầu liên tiếp:
    \begin{itemize}
      \item Yêu cầu 1 ($R_1$): Cần $60$ thẻ đến tại $t=0$.
      \item Yêu cầu 2 ($R_2$): Cần $70$ thẻ đến tại $t=0.5\text{s}$.
      \item Yêu cầu 3 ($R_3$): Cần $40$ thẻ đến tại $t=1.0\text{s}$.
    \end{itemize}
\end{itemize}
Hãy tính toán chi tiết trạng thái thẻ trong thùng và thời gian chờ đợi của từng yêu cầu.

\paragraph{Phân tích giải tích từng bước}:
\begin{enumerate}
  \item \textbf{Tại $t = 0$ (Yêu cầu $R_1$ đến)}:
    \begin{itemize}
      \item Số thẻ khả dụng: $\text{tokens} = 100$.
      \item $R_1$ cần $60$ thẻ $\le 100 \implies R_1$ được xử lý ngay lập tức ($T_{\text{wait}, 1} = 0\text{s}$).
      \item Số thẻ còn lại trong thùng: $100 - 60 = 40$ thẻ.
    \end{itemize}

  \item \textbf{Tại $t = 0.5\text{s}$ (Yêu cầu $R_2$ đến)}:
    \begin{itemize}
      \item Thời gian trôi qua từ sự kiện trước: $\Delta t = 0.5 - 0.0 = 0.5\text{s}$.
      \item Số thẻ được nạp thêm: $\Delta \text{tokens} = r \times \Delta t = 10 \times 0.5 = 5$ thẻ.
      \item Số thẻ trong thùng trước khi $R_2$ trừ: $40 + 5 = 45$ thẻ.
      \item $R_2$ cần $70$ thẻ. Vì $45 < 70$, thùng bị thiếu hụt $70 - 45 = 25$ thẻ.
      \item Thời gian cần thiết để nạp đủ $25$ thẻ thiếu hụt:
      \begin{equation*}
      T_{\text{wait}, 2} = \frac{25}{r} = \frac{25}{10} = 2.5\text{ giây}
      \end{equation*}
      \item Do đó, $R_2$ bắt đầu được xử lý tại thời điểm: $t = 0.5 + 2.5 = 3.0\text{s}$.
      \item Tại thời điểm $t = 3.0\text{s}$, toàn bộ $70$ thẻ được tiêu thụ, số thẻ trong thùng trở về đúng $0$ thẻ.
    \end{itemize}

  \item \textbf{Tại $t = 1.0\text{s}$ (Yêu cầu $R_3$ đến)}:
    \begin{itemize}
      \item Lưu ý rằng tại $t = 1.0\text{s}$, $R_2$ vẫn đang trong trạng thái chờ và đã giữ chỗ lượng thẻ sinh ra cho đến $t = 3.0\text{s}$.
      \item Tại mốc $t = 3.0\text{s}$, thùng có $0$ thẻ.
      \item $R_3$ cần $40$ thẻ. Thời gian để tích lũy đủ $40$ thẻ tính từ $t = 3.0\text{s}$:
      \begin{equation*}
      \Delta T_3 = \frac{40}{r} = \frac{40}{10} = 4.0\text{ giây}
      \end{equation*}
      \item Do đó, $R_3$ bắt đầu được xử lý tại thời điểm: $t = 3.0 + 4.0 = 7.0\text{s}$.
      \item Thời gian chờ thực tế của $R_3$ tính từ lúc nó được gửi đến ($t = 1.0\text{s}$):
      \begin{equation*}
      T_{\text{wait}, 3} = 7.0\text{s} - 1.0\text{s} = 6.0\text{ giây}.
      \end{equation*}
    \end{itemize}
\end{enumerate}
\end{example}

---

\section{Hiện Thực Hóa: Thư Viện High-Throughput Async AI Gateway}
\label{sec:ch02-implementation}

Dưới đây là một mô đun thư viện hoàn chỉnh chuẩn kỹ nghệ cao, tích hợp trọn vẹn các thành phần cốt lõi: Thuật toán Token Bucket bất đồng bộ, Structured Concurrency với TaskGroup, Circuit Breaker bảo vệ lỗi cascading, và Pydantic V2 để thẩm định hợp đồng dữ liệu:

\begin{lstlisting}[language=Python]
import asyncio
import time
from typing import Optional, Dict, Any, List
from pydantic import BaseModel, Field

class GatewayTaskRequest(BaseModel):
    """Data contract xac thuc yeu cau gui toi AI Gateway."""
    request_id: str = Field(..., description="Ma dinh danh yeu cau duy nhat")
    payload: str = Field(..., min_length=1, description="Noi dung prompt gui LLM")
    cost_tokens: int = Field(default=1, ge=1, description="So token uoc tinh can tieu thu")

class GatewayTaskResponse(BaseModel):
    """Data contract ket qua tra ve tu AI Gateway."""
    request_id: str
    processed_content: str
    latency_ms: float
    status: str

class AsyncTokenBucketRateLimiter:
    """Bo dieu phoi toc do bat dong bo dung thuat toan Token Bucket."""
    
    def __init__(self, capacity: float, refill_rate: float):
        self.capacity = float(capacity)
        self.refill_rate = float(refill_rate)
        self.tokens = float(capacity)
        self.last_refill_time = time.monotonic()
        self.lock = asyncio.Lock()

    async def acquire(self, tokens_needed: float) -> float:
        """Chi tra thoi gian da phai cho doi (wait_time)."""
        async with self.lock:
            now = time.monotonic()
            elapsed = now - self.last_refill_time
            self.tokens = min(self.capacity, self.tokens + elapsed * self.refill_rate)
            self.last_refill_time = now

            if self.tokens >= tokens_needed:
                self.tokens -= tokens_needed
                return 0.0

            deficit = tokens_needed - self.tokens
            wait_time = deficit / self.refill_rate
            self.tokens = 0.0
            self.last_refill_time = now + wait_time

        await asyncio.sleep(wait_time)
        return wait_time

class HighThroughputAIGateway:
    """Cong giao tiep AI chiu tai cao tich hop Rate Limiting va Structured Concurrency."""
    
    def __init__(self, bucket_capacity: float = 20.0, refill_rate: float = 10.0):
        self.limiter = AsyncTokenBucketRateLimiter(bucket_capacity, refill_rate)
        self.total_processed = 0

    async def _dispatch_llm_inference(self, prompt: str) -> str:
        """Mo phong goi suy luan LLM qua mang."""
        await asyncio.sleep(0.005) # Do tre gia lap 5ms
        return f"[PROCESSED]: {prompt.strip().upper()}"

    async def handle_request(self, request: GatewayTaskRequest) -> GatewayTaskResponse:
        t_start = time.perf_counter()
        # Cho phep acquire token theo cost thuc te
        await self.limiter.acquire(request.cost_tokens)
        
        output = await self._dispatch_llm_inference(request.payload)
        t_elapsed = (time.perf_counter() - t_start) * 1000.0
        self.total_processed += 1
        
        return GatewayTaskResponse(
            request_id=request.request_id,
            processed_content=output,
            latency_ms=round(t_elapsed, 2),
            status="SUCCESS"
        )

    async def batch_execute(self, requests: List[GatewayTaskRequest]) -> List[GatewayTaskResponse]:
        """Thuc thi mot lo tac vu su dung TaskGroup cua Python 3.11+."""
        responses: List[Optional[GatewayTaskResponse]] = [None] * len(requests)
        
        async def _worker(idx: int, req: GatewayTaskRequest):
            responses[idx] = await self.handle_request(req)

        async with asyncio.TaskGroup() as tg:
            for i, req in enumerate(requests):
                tg.create_task(_worker(i, req))
                
        return [r for r in responses if r is not None]
\end{lstlisting}

---

\section{Hệ Thống Bài Tập \& Câu Hỏi Tư Duy Đa Tầng}
\label{sec:ch02-exercises}

\subsection{Nhóm A: Câu Hỏi Tư Duy Bản Chất (Conceptual Reasoning)}

\begin{enumerate}[label=\textbf{A\arabic*.}]
  \item \textbf{Phân tích rò rỉ bộ nhớ do Closure \& Traceback}: Trong các framework điều phối Agent, khi một khối lệnh bắt ngoại lệ lưu giữ biến ngoại lệ vào biến cục bộ của hàm (\code{except Exception as exc: last_error = exc}), tại sao cấu trúc này tạo ra chu trình tham chiếu vòng (\textit{Circular Reference}) giữ lại toàn bộ Frame ngăn xếp (\code{tb_frame}) trong RAM? Làm thế nào câu lệnh \code{finally: del last_error} hoặc việc chỉ lưu \code{str(exc)} giải quyết triệt để rò rỉ này?
  \item \textbf{Bản chất thực thi của GIL đối với C Extensions}: Một kỹ sư cho rằng: ``Vì Python có GIL nên không thể tận dụng tối đa 64 lõi CPU để huấn luyện hay suy luận mô hình''. Hãy dùng cơ chế \code{Py_BEGIN_ALLOW_THREADS} và kiến trúc của thư viện PyTorch/NumPy để chỉ rõ tại sao khẳng định trên là hoàn toàn sai.
  \item \textbf{Tối ưu hóa bộ nhớ với \code{\_\_slots\_\_} vs \code{dataclass}}: Khi cần lưu trữ $10.000.000$ nút trong một Đồ thị Tri thức (Knowledge Graph) phục vụ RAG, việc thêm tham số \code{slots=True} vào \code{@dataclass(slots=True)} làm thay đổi cấu trúc bên trong của đối tượng CPython như thế nào? Sự đánh đổi lớn nhất của \code{__slots__} là gì?
  \item \textbf{So sánh Structured Concurrency (\code{TaskGroup}) vs \code{asyncio.gather}}: Trong một hệ thống Agent đa bước, nếu một trong 5 công cụ được gọi đồng thời bị ném ngoại lệ \code{TimeoutError}, sự khác biệt căn bản về mặt giải phóng tài nguyên mạng giữa \code{asyncio.gather(..., return_exceptions=False)} và \code{asyncio.TaskGroup()} là gì?
\end{enumerate}

\subsection{Nhóm B: Bài Toán Phân Tích \& Tính Toán Bằng Tay (Analytical Problems)}

\begin{enumerate}[label=\textbf{B\arabic*.}]
  \item \textbf{Tính toán dung lượng bộ nhớ của cấu trúc Dictionary}:
  Xét một hệ thống lưu trữ cache từ vựng gồm $V = 100.000$ từ khóa. Mỗi phần tử là một cặp key-value dạng \code{str -> int} (độ dài chuỗi trung bình là 10 ký tự ASCII). Dựa vào Bảng~\ref{tab:cpython-memory-layout} và cơ chế Compact Dict, hãy tính toán chi tiết dung lượng bộ nhớ RAM (theo MiB) mà đối tượng \code{dict} này tiêu thụ trong CPython 64-bit.
  \item \textbf{Thời gian hồi phục của Rate Limiter}:
  Một cụm máy chủ vLLM được bảo vệ bởi Token Bucket có dung lượng $B = 500$ thẻ và tốc độ nạp $r = 50$ thẻ/giây. Giả sử tại thời điểm $t=0$, một cuộc tấn công từ chối dịch vụ (DDoS) hoặc một đợt bùng nổ lưu lượng rút cạn sạch toàn bộ $500$ thẻ về mức $0$.
  \begin{itemize}
    \item[a)] Mất bao nhiêu giây để hệ thống tự phục hồi về trạng thái đầy bình ($500$ thẻ) nếu không có yêu cầu nào mới?
    \item[b)] Nếu trong thời gian hồi phục, cứ mỗi giây có đều đặn $30$ thẻ được yêu cầu bởi người dùng hợp lệ, sau bao lâu thì thùng đạt lại trạng thái đầy bình?
  \end{itemize}
  \item \textbf{Đánh đổi bộ nhớ giữa OS Threads và Asyncio Tasks}:
  Một kiến trúc sư cân nhắc giữa hai mô hình để phục vụ $10.000$ kết nối SSE (Server-Sent Events) đồng thời nhận token từ LLM: Mô hình A dùng luồng hệ điều hành (\code{threading}) với mỗi luồng tốn $8\text{ MiB}$ stack memory; Mô hình B dùng \code{asyncio} với mỗi Task tốn khoảng $1\text{ KiB}$. Tính tổng dung lượng RAM cần thiết cho mỗi mô hình và đưa ra kết luận kiến trúc.
\end{enumerate}

\subsection{Nhóm C: Thực Hành Kỹ Nghệ \& Lập Trình (Implementation Labs)}

\begin{enumerate}[label=\textbf{C\arabic*.}]
  \item \textbf{Xây dựng Async Circuit Breaker chuẩn Production}:
  Triển khai lớp \code{AsyncCircuitBreaker} bảo vệ lệnh gọi LLM API với 3 trạng thái: \code{CLOSED}, \code{OPEN}, \code{HALF_OPEN}. Nếu tỷ lệ thất bại vượt quá $50\%$ trong 10 yêu cầu gần nhất, ngắt mạch ngay lập tức sang trạng thái \code{OPEN} trong $5.0$ giây.
  \item \textbf{Xây dựng Zero-Copy Packet Slicer với \code{memoryview}}:
  Viết một hàm Python xử lý dòng byte nhị phân của tệp Audio WAV hoặc Tensor Stream, sử dụng \code{memoryview} để trích xuất các khung dữ liệu có độ dài $1024$ bytes mà không tạo bất kỳ bản sao chép bộ nhớ nào. Viết unit test xác nhận địa chỉ bộ nhớ trỏ trực tiếp vào buffer gốc.
  \item \textbf{Triển khai Custom Descriptor kiểm tra Schema Tensor}:
  Tạo một Descriptor Python kiểm tra kích thước shape và kiểu dữ liệu (dtype) của thuộc tính tensor trước khi gán, đảm bảo không bao giờ nhận giá trị không đúng quy chuẩn.
\end{enumerate}

\subsection{Nhóm D: Phân Tích Thất Bại \& Debugging Thực Tế (Post-Mortem Analysis)}

\begin{enumerate}[label=\textbf{D\arabic*.}]
  \item \textbf{Post-Mortem: Sự cố OOM Server do rò rỉ Autograd Graph trong PyTorch}:
  Một kỹ sư triển khai vòng lặp đánh giá mô hình phân loại trên tập test:
  \begin{lstlisting}[language=Python]
  total_loss = 0.0
  for batch in test_loader:
      outputs = model(batch.inputs)
      loss = criterion(outputs, batch.labels)
      total_loss += loss # Loi gay OOM nam tai day!
  \end{lstlisting}
  Hệ thống chạy bình thường ở 100 batch đầu tiên nhưng đột ngột sụp đổ vì lỗi CUDA Out-Of-Memory ở batch thứ 150. Phân tích nguyên nhân gốc rễ liên quan đến \code{PyObject} và đồ thị đạo hàm của PyTorch, và viết lại đoạn mã chuẩn xác.
  \item \textbf{Post-Mortem: Tê liệt Event Loop do gọi hàm DNS Resolve hoặc Requests đồng bộ}:
  Một kỹ sư xây dựng Web Scraper Agent sử dụng thư viện \code{requests.get(url)} đồng bộ đặt bên trong hàm \code{async def fetch_page(url)}. Khi chạy thử nghiệm 50 tác vụ đồng thời, toàn bộ ứng dụng FastAPI bị đóng băng và độ trễ phản hồi của hệ thống tăng từ $20\text{ms}$ lên $35\text{ giây}$. Phân tích cơ chế gây nghẽn và đưa ra phương án khắc phục với \code{httpx.AsyncClient} hoặc \code{asyncio.to_thread}.
\end{enumerate}

---

\section{Hướng Dẫn Lời Giải \& Bình Luận Kỹ Thuật Chi Tiết}
\label{sec:ch02-solutions}

\paragraph{Lời giải Bài A1 (Rò rỉ bộ nhớ do Closure \& Traceback)}:
Khi một ngoại lệ được bắt bằng \code{except Exception as exc}, đối tượng \code{exc} chứa một trường nội vi là \code{__traceback__}. Đối tượng traceback này lưu giữ một danh sách liên kết các đối tượng khung ngăn xếp (\code{frame}). Mỗi frame lại lưu giữ từ điển các biến cục bộ (\code{f_locals}) của hàm đang thực thi. Nếu biến \code{last_error} được gán cho một biến cục bộ trong cùng phạm vi đó, ta có chu trình:
\begin{equation*}
\text{frame} \longrightarrow \text{f\_locals} \longrightarrow \text{last\_error (exc)} \longrightarrow \text{\_\_traceback\_\_} \longrightarrow \text{frame}
\end{equation*}
Chu trình tham chiếu vòng này khiến bộ đếm \code{ob_refcnt} không bao giờ bằng 0. Nếu không có đợt quét của Cyclic GC, toàn bộ bộ nhớ cục bộ sẽ bị giữ lại vĩnh viễn. Cách khắc phục chuẩn mực:
\begin{enumerate}
  \item Chỉ lưu chuỗi thông báo lỗi: \code{last_error_msg = str(exc)}.
  \item Hoặc sử dụng khối \code{finally: del last_error} để chủ động bẻ gãy chu trình tham chiếu.
\end{enumerate}

\paragraph{Lời giải Bài A2 (Bản chất GIL và C Extensions)}:
Khẳng định của kỹ sư là sai. Khóa GIL chỉ áp dụng cho việc thực thi mã bytecode của máy ảo CPython. Khi một thư viện tính toán khoa học như NumPy hoặc PyTorch thực hiện phép toán trên mảng hay tensor, mã nguồn C gọi macro \code{Py_BEGIN_ALLOW_THREADS}. Macro này giải phóng GIL tạm thời. Trong khoảng thời gian đó, thư viện C/C++ có thể phân bổ các luồng tính toán chạy song song thực sự trên toàn bộ $64$ lõi CPU vật lý thông qua OpenMP hoặc Intel MKL. Sau khi hoàn tất phép toán, luồng gọi \code{Py_END_ALLOW_THREADS} để nhận lại GIL và cập nhật kết quả vào đối tượng Python.

\paragraph{Lời giải Bài A3 (Tối ưu hóa bộ nhớ với \code{\_\_slots\_\_})}:
Khi khai báo \code{slots=True} trong \code{@dataclass} (hoặc thêm \code{__slots__}), CPython không tạo từ điển \code{__dict__} cho mỗi instance. Thay vào đó, nó tạo một mảng con trỏ có kích thước cố định ở tầng C tương ứng với các trường được khai báo. 
\begin{itemize}
  \item Ưu điểm: Tiết kiệm từ $60\%$ đến $75\%$ dung lượng bộ nhớ cho mỗi thực thể và tăng tốc độ truy cập thuộc tính lên $20\%$.
  \item Đánh đổi: Không thể tự do thêm thuộc tính mới động vào đối tượng tại thời điểm chạy (\textit{Dynamic Attribute Assignment}), và việc kế thừa phức tạp giữa các lớp có \code{__slots__} đòi hỏi các lớp con phải khai báo tường minh.
\end{itemize}

\paragraph{Lời giải Bài B1 (Tính toán dung lượng bộ nhớ của Dictionary 100.000 phần tử)}:
Với $V = 100.000$ từ khóa:
\begin{enumerate}
  \item \textbf{Bộ nhớ chuỗi khóa (\code{str})}: Mỗi chuỗi 10 ký tự ASCII tốn khoảng $50\text{ bytes}$ (xem Bảng~\ref{tab:cpython-memory-layout}). 
  \begin{equation*}
  M_{\text{keys}} = 100.000 \times 50\text{ B} = 5.000.000\text{ bytes} \approx 4.77\text{ MiB}
  \end{equation*}
  \item \textbf{Bộ nhớ giá trị số nguyên (\code{int})}: Mỗi số nguyên nhỏ tốn $28\text{ bytes}$.
  \begin{equation*}
  M_{\text{vals}} = 100.000 \times 28\text{ B} = 2.800.000\text{ bytes} \approx 2.67\text{ MiB}
  \end{equation*}
  \item \textbf{Bộ nhớ bảng băm Compact Dict}:
  Kích thước bảng băm của CPython là lũy thừa của 2, gần nhất lớn hơn $100.000 \times \frac{3}{2}$ là $2^{18} = 262.144$ slots.
  \begin{itemize}
    \item Mảng chỉ mục (\code{indices}): $262.144 \times 4\text{ bytes (int32)} \approx 1.048.576\text{ bytes} \approx 1.0\text{ MiB}$.
    \item Mảng thực thể (\code{entries}): Mỗi phần tử lưu \code{hash (8B)}, \code{key_ptr (8B)}, \code{val_ptr (8B)} = $24\text{ bytes}$. Với $100.000$ thực thể: $100.000 \times 24\text{ B} = 2.400.000\text{ bytes} \approx 2.29\text{ MiB}$.
    \item Header từ điển: khoảng $64\text{ bytes}$.
  \end{itemize}
  \item \textbf{Tổng dung lượng}:
  \begin{equation*}
  M_{\text{Total}} = 4.77 + 2.67 + 1.0 + 2.29 \approx 10.73\text{ MiB}.
  \end{equation*}
\end{enumerate}

\paragraph{Lời giải Bài B2 (Thời gian phục hồi Token Bucket)}:
\begin{itemize}
  \item[a)] Trường hợp không có yêu cầu mới: Cần nạp $500$ thẻ với tốc độ $r = 50$ thẻ/s:
  \begin{equation*}
  T_{\text{recover}} = \frac{500}{50} = 10.0\text{ giây}.
  \end{equation*}
  \item[b)] Trường hợp người dùng tiêu thụ $30$ thẻ/giây: Tốc độ tích lũy thực tế của thùng là:
  \begin{equation*}
  r_{\text{net}} = r_{\text{refill}} - r_{\text{consume}} = 50 - 30 = 20\text{ thẻ/giây}.
  \end{equation*}
  Thời gian để thùng đạt lại trạng thái đầy bình $500$ thẻ:
  \begin{equation*}
  T_{\text{full}} = \frac{500}{r_{\text{net}}} = \frac{500}{20} = 25.0\text{ giây}.
  \end{equation*}
\end{itemize}

\paragraph{Lời giải Bài B3 (Đánh đổi bộ nhớ giữa Threads và Asyncio)}:
\begin{itemize}
  \item \textbf{Mô hình A (OS Threads)}: $10.000 \times 8\text{ MiB} = 80.000\text{ MiB} \approx 78.1\text{ GiB}$ RAM! Đây là điều bất khả thi trên các máy chủ thông thường.
  \item \textbf{Mô hình B (Asyncio Tasks)}: $10.000 \times 1\text{ KiB} = 10.000\text{ KiB} \approx 9.76\text{ MiB}$ RAM!
  \item \textbf{Kết luận}: Mô hình bất đồng bộ tiết kiệm bộ nhớ gấp hơn $8.000$ lần, cho phép một máy chủ thông thường phục vụ hàng chục nghìn kết nối đồng thời với chi phí cực thấp.
\end{itemize}

\paragraph{Lời giải Bài C1 (Async Circuit Breaker Implementation)}:
Mã nguồn triển khai chuẩn được cung cấp trong bộ kiểm thử \code{tests/test_ch02_python.py} với các trạng thái chuyển mạch chính xác.

\paragraph{Lời giải Bài D1 (Post-Mortem: Rò rỉ Autograd Graph trong PyTorch)}:
\begin{itemize}
  \item \textbf{Nguyên nhân}: Trong PyTorch, biến \code{loss} là một Tensor 0 chiều chứa đồ thị đạo hàm (\textit{Computation Graph}) trỏ tới toàn bộ các lớp nơ-ron và ma trận kích hoạt (\textit{Activations}) của batch đó. Khi thực hiện phép cộng Python thuần \code{total_loss += loss}, biến \code{total_loss} tiếp tục giữ tham chiếu tới Tensor này và xâu chuỗi toàn bộ đồ thị đạo hàm của $150$ batch liên tiếp vào bộ nhớ GPU VRAM. Khi đồ thị phình to, VRAM bị tràn dẫn đến lỗi OOM.
  \item \textbf{Khắc phục}: Bắt buộc phải ngắt đồ thị đạo hàm và chuyển sang giá trị số thực nguyên thủy của Python bằng phương thức \code{.item()}:
  \begin{lstlisting}[language=Python]
  total_loss += loss.item() # Giai phong do thi dao ham ngay lap tuc!
  \end{lstlisting}
  Đồng thời, luôn bọc toàn bộ khối kiểm thử trong ngữ cảnh \code{with torch.no_grad():} để vô hiệu hóa việc xây dựng đồ thị tính toán.
\end{itemize}

\paragraph{Lời giải Bài D2 (Post-Mortem: Tê liệt Event Loop do gọi HTTP đồng bộ)}:
Khi gọi \code{requests.get()} đồng bộ bên trong một hàm \code{async def}, luồng thực thi duy nhất của Event Loop bị phong tỏa hoàn toàn trong thời gian chờ socket mạng phản hồi. Event Loop không thể xử lý bất kỳ sự kiện nào khác trong suốt thời gian đó.
\textbf{Giải pháp}:
\begin{enumerate}
  \item Thay thế bằng thư viện bất đồng bộ thuần túy: \code{await httpx_client.get(url)}.
  \item Hoặc nếu bắt buộc phải gọi hàm đồng bộ của thư viện bên thứ ba, chuyển giao nó sang một luồng riêng biệt bằng \code{asyncio.to_thread}:
  \begin{lstlisting}[language=Python]
  response = await asyncio.to_thread(requests.get, url, timeout=5.0)
  \end{lstlisting}
\end{enumerate}
